% Kapitel 2: Problemdefinition & Use Case
\section{Problemdefinition \& Use Case}

\subsection{Nutzungssituation}

Die Website des Immersive Reality Lab richtet sich an mehrere heterogene Nutzergruppen:

\begin{itemize}
    \item \textbf{Studierende}, die nach Bachelor-/Masterarbeitsthemen, Lehrveranstaltungen oder Geräteausleihe suchen.
    \item \textbf{Forschungsinteressierte}, die Publikationen und Kooperationsmöglichkeiten erkunden möchten.
    \item \textbf{Industriepartner}, die sich über konkrete Projekte, Technologietransfer und Ansprechpartner informieren wollen.
\end{itemize}

Die Nutzungssituation ist typischerweise \textit{goal-oriented}: Der Besucher hat eine konkrete Frage und möchte diese schnell beantwortet haben, ohne die gesamte Seitenstruktur durchsuchen zu müssen.

\subsection{Defizite der aktuellen Website}

Eine Analyse der bestehenden Website zeigt folgende Problembereiche:

\begin{enumerate}
    \item \textbf{heterogene Inhaltsstruktur}: Projekte, Publikationen und Lehre sind über verschiedene Unterseiten verteilt, teils mit dynamischen Parametern.
    \item \textbf{Keine interne Suchfunktion}: Nutzer müssen manuell navigieren oder die Browsersuche (Strg+F) verwenden.
    \item \textbf{Fachspezifische Terminologie}: Begriffe wie ``INFUSE'', ``XRT-HuFuSa'' oder ``Biosensorik'' setzen Vorwissen voraus.
\end{enumerate}

\subsection{Anforderungen an den Chatbot}

Aus der Nutzungssituation leiten sich folgende Anforderungen ab:

\begin{tabularx}{\textwidth}{lX}
\toprule
\textbf{Kategorie} & \textbf{Anforderung} \\
\midrule
Funktional & Beantwortung von Fragen zu allen Website-Inhalten \\
Funktional & Quellenangabe für jede Antwort \\
Funktional & Gesprächsverlauf über mehrere Anfragen \\
Usability & Onboarding: Chatbot erklärt eigene Fähigkeiten \\
Usability & Fehlertoleranz: klare Rückmeldung bei Wissenslücken \\
Usability & Ladeindikator während der Antwortgenerierung \\
Technisch & Datenschutz: lokale Verarbeitung (Ollama, ChromaDB) \\
\bottomrule
\end{tabularx}
